\section{Introduction}
\label{sec:intro}

As language models become increasingly capable, they are being deployed in increasingly complex applications. In such contexts, the assumption of a single ground truth label often harms the performance of the model, and increases the risk of harm to users \citep{plank-2022-problem,pmlr-v235-sorensen24a}. Much of the recent work on improving pluralistic reasoning in language models has focused on \textit{generation} by considering approaches like reinforcement learning from human feedback (RLHF) \citep{ouyang2022training} or from AI feedback (RLAIF) \citep{bai2022constitutionalaiharmlessnessai}. However, many applications of language models involve \textit{retrieval} and \textit{synthesis} of existing information, as well as generation, and these earlier steps are often overlooked in discussions of pluralism. 

In this paper, we argue that pluralism is important not only in the generation step, but also earlier in the process, and that it is important to consider all steps when designing retrieval-augmented generation (RAG) pipelines. We also discuss some of the challenges and opportunities for improving pluralistic reasoning in RAG pipelines. As \citet{agrawal2026retrievalaugmentedgenerationfactualgrounding} argue, RAG pipelines exhibit a factual bias which reduces the diversity of information retrieved and synthesized. While opinion-aware capabilities are important, they are not sufficient to ensure pluralistic reasoning in RAG pipelines. We argue that to ensure pluralistic reasoning while still providing accurate and relevant information, RAG pipelines must be sensitive to principles of \textit{social choice} \citep{arrow1951socialchoice} at the retrieval and synthesis stages.
